\chapter{Automated Market Makers}\label{ch:m3:automated-market-makers}

A pool holds ether and dollars and quotes no prices. Anyone may send it one \omterm{def:m3:blockchains-for-traders:key}{token} and take out the
other, in any amount, provided that after the trade the product of its two balances is no smaller than
before. That single rule, written in a few dozen lines of a \omterm{def:m3:blockchains-for-traders:contract}{smart contract}, makes a market: the price is
the ratio of the balances, a buyer pushes it up by taking ether out, and arbitrageurs keep it in line
with the price on \omterm{def:m3:centralised-exchanges:cex}{centralised exchanges} by trading whenever the two differ. Those who deposit the
balances earn the fees and bear a cost that is less obvious than the fees and, for most pools, larger.
This chapter explains the \omterm{def:m3:automated-market-makers:amm}{automated market maker}: the constant-product rule and its concentrated
variant, the invariant used for \omterm{def:m3:blockchains-for-traders:stable}{stablecoins}, what liquidity providers earn and lose, measured the way
the recent literature measures it, and the routers and auctions that sit between traders and pools.

\section{Constant-product market makers}

\begin{definition}[Decentralised exchange, automated market maker, liquidity pool]\label{def:m3:automated-market-makers:amm}
A \emph{decentralised exchange}\index{decentralised exchange} is a trading venue implemented as \omterm{def:m3:blockchains-for-traders:contract}{smart
contracts} on a \omterm{def:m3:blockchains-for-traders:chain}{blockchain}, on which users trade from their own wallets and every trade settles on chain.
An \emph{automated market maker}\index{automated market maker} is a decentralised exchange whose prices are
set by a formula applied to the balances of a \emph{liquidity pool}\index{liquidity pool}: a contract holding
reserves of two or more \omterm{def:m3:blockchains-for-traders:key}{tokens}, deposited by liquidity providers who receive a claim on them and on the fees.
\end{definition}

\begin{definition}[Constant-product market maker]\label{def:m3:automated-market-makers:cp}
A \emph{constant-product market maker}\index{constant-product market maker} is an \omterm{def:m3:automated-market-makers:amm}{automated market maker} whose
pool of reserves $x$ and $y$ accepts a trade only if $xy$ does not decrease, with a fee $\gamma_{\mathrm{fee}}$ taken on
the input and left in the pool.
\end{definition}

\begin{proposition}[Swap output and marginal price]\label{prop:m3:automated-market-makers:swap}
A trader who sends $\Delta x$ to a constant-product pool with reserves $(x, y)$ and fee $\gamma_{\mathrm{fee}}$ receives
\[
\Delta y = \frac{y\,(1 - \gamma_{\mathrm{fee}})\,\Delta x}{x + (1 - \gamma_{\mathrm{fee}})\,\Delta x}.
\]
The pool's marginal price of the first \omterm{def:m3:blockchains-for-traders:key}{token} is $y/x$; the average price of the trade is $\Delta y / \Delta x$, below
$y/x$ by the fee and by a price impact that grows with $\Delta x / x$. Without the fee, the product is exactly preserved;
with it, the product grows.
\end{proposition}

\begin{proof}
Only $(1 - \gamma_{\mathrm{fee}})\Delta x$ counts towards the invariant: $(x + (1 - \gamma_{\mathrm{fee}})\Delta x)(y - \Delta y) = xy$
gives $\Delta y$. The fee part of $\Delta x$ is added to the reserves on top, so $(x + \Delta x)(y - \Delta y) > xy$.
The marginal price is $-dy/dx$ along $xy = k$, equal to $y/x$.
\end{proof}

A pool of 1\,000 ether and 3 million dollars quotes ether at 3\,000 dollars. Selling it 10 ether with a 0.30\% fee returns
29\,614.74 dollars, an average of 2\,961.47: 128 basis points below 3\,000, of which 30 are the fee and about 98 the impact. The pool has no bid and no
offer; its ``spread'' is twice the fee plus the impact of the size, and its depth is set by its reserves. Its price moves
only when someone trades with it. When ether moves on \omterm{def:m3:centralised-exchanges:cex}{centralised exchanges}, the pool's price is stale until an
arbitrageur trades it back into line, and the arbitrageur's profit comes out of the pool.

\begin{dated}{2026-09}{Fee tiers and versions of the largest AMM}\label{dat:m3:automated-market-makers:uniswap}
Uniswap v2 charges a fee of 30 basis points on trades, paid to liquidity providers. Uniswap v3 (whitepaper dated March 2021, deployed to
Ethereum mainnet on 5 May 2021)
introduced \omterm{def:m3:automated-market-makers:concentrated}{concentrated liquidity} and fee tiers, initially 0.05\%, 0.30\% and 1\% with tick spacings of 10, 60 and 200 ticks,
where a tick is a price step of a factor 1.0001. Uniswap v4 became operational on 31 January 2025, adding hooks, modular plugins
that let developers add custom logic for pools, swaps, fees and positions. By CoinGecko's count, spot volume on \omterm{def:m3:automated-market-makers:amm}{decentralised exchanges}
rose from 6.0\% of \omterm{def:m3:centralised-exchanges:cex}{centralised exchanges}' spot volume in January 2021 to a peak of 37.4\% in June 2025, and stood near 21\% in November 2025.
\end{dated}

\section{Concentrated liquidity}

\begin{definition}[Concentrated liquidity]\label{def:m3:automated-market-makers:concentrated}
\emph{Concentrated liquidity}\index{concentrated liquidity} is liquidity provided to an \omterm{def:m3:automated-market-makers:amm}{automated market maker} over a chosen
price range only: the position behaves as a constant-product pool with a larger virtual reserve while the price is inside
the range, and holds only one of the two \omterm{def:m3:blockchains-for-traders:key}{tokens} when the price is outside it.
\end{definition}

In Uniswap v3 a position is a liquidity amount $L$ over a range $[P_a, P_b)$. Inside the range, with square-root price
$\sqrt{P}$, it holds $L(1/\sqrt{P} - 1/\sqrt{P_b})$ of the first \omterm{def:m3:blockchains-for-traders:key}{token} and $L(\sqrt{P} - \sqrt{P_a})$ of the second; below
$P_a$ it holds only the first \omterm{def:m3:blockchains-for-traders:key}{token}, above $P_b$ only the second. Prices are stored as square roots in binary fixed point,
and ranges end on ticks. A narrow range needs far fewer \omterm{def:m3:blockchains-for-traders:key}{tokens} than a full-range position for the same $L$, so the same
capital provides more depth near the current price: a range of $\pm10\%$ is about 21 times as capital-efficient as the
full range. It also sells its \omterm{def:m3:blockchains-for-traders:key}{tokens} faster as the price moves through the range, and it earns nothing once the price has
left it.

\begin{omfigure}
\begin{tikzpicture}
\begin{axis}[width=11cm,height=5cm,xmode=log,log basis x=2,
  xlabel={price relative to entry},ylabel={value against holding, \%},
  tick label style={font=\small},label style={font=\small},
  xmin=0.5,xmax=2,ymin=-20,ymax=0.5,grid=major,grid style={gray!25},
  xtick={0.5,0.707,1,1.414,2},xticklabels={0.5,0.71,1,1.41,2},
  legend columns=3,legend style={font=\footnotesize,at={(0.5,-0.3)},anchor=north,draw=none,fill=none,/tikz/every even column/.append style={column sep=8pt}},
  table/col sep=comma]
\addplot[omThm,very thick] table[x=r,y=full]{figdata/markets-3/20-automated-market-makers/il.csv};\addlegendentry{full range}
\addplot[omDef,very thick,dashed] table[x=r,y=w50]{figdata/markets-3/20-automated-market-makers/il.csv};\addlegendentry{range $\pm50\%$}
\addplot[omExo,very thick,dotted] table[x=r,y=w10]{figdata/markets-3/20-automated-market-makers/il.csv};\addlegendentry{range $\pm10\%$}
\end{axis}
\end{tikzpicture}

\omcaption{\omterm{def:m3:automated-market-makers:il}{Impermanent loss}: value of a liquidity position against simply holding the \omterm{def:m3:blockchains-for-traders:key}{tokens} deposited, after the price moves by a
factor, for a full-range constant-product position and for concentrated positions over $\pm50\%$ and $\pm10\%$ around the entry price.
The narrower the range, the faster the loss accrues; outside the range the position holds only the \omterm{def:m3:blockchains-for-traders:key}{token} that has lost value
against the other. Data: the chapter's tutorial.}\label{fig:m3:automated-market-makers:il}
\end{omfigure}

\section{Stableswap}

For two \omterm{def:m3:blockchains-for-traders:stable}{stablecoins} that should trade at one for one, the constant product wastes liquidity: most of it sits at prices far from one
that should never trade.

\begin{definition}[Stableswap]\label{def:m3:automated-market-makers:stableswap}
A \emph{stableswap}\index{stableswap} market maker is an \omterm{def:m3:automated-market-makers:amm}{automated market maker} whose invariant behaves like a constant sum
(no price impact) near balanced reserves and like a constant product far from balance, blended by an amplification coefficient $A$:
for $n$ coins with balances $x_i$ and invariant $D$,
\[
A n^n \sum_i x_i + D = A D n^n + \frac{D^{n+1}}{n^n \prod_i x_i}.
\]
\end{definition}

The invariant is solved numerically, for $D$ when liquidity is added and for one balance when a trade is made, by integer Newton
iterations. With $A$ large the pool trades near one for one over a wide range of balances; if one coin \omterm{def:m3:blockchains-for-traders:stable}{depegs}, arbitrageurs sell it into
the pool until the pool holds mostly that coin, and the price then moves as in a constant-product pool. The liquidity providers of a
\omterm{def:m3:blockchains-for-traders:stable}{stablecoin} pool are therefore sellers of a put on the weaker coin's peg.

\begin{omfigure}
\begin{tikzpicture}
\begin{axis}[width=11cm,height=5cm,
  xlabel={balance of coin X},ylabel={balance of coin Y},
  tick label style={font=\small},label style={font=\small},
  xmin=0,xmax=20,ymin=0,ymax=20,grid=major,grid style={gray!25},
  legend columns=3,legend style={font=\footnotesize,at={(0.5,-0.3)},anchor=north,draw=none,fill=none,/tikz/every even column/.append style={column sep=8pt}},
  table/col sep=comma]
\addplot[omThm,very thick] table[x=x,y=cp]{figdata/markets-3/20-automated-market-makers/invariants.csv};\addlegendentry{constant product}
\addplot[gray,thick,dashed,restrict x to domain=0:10] table[x=x,y=csum]{figdata/markets-3/20-automated-market-makers/invariants.csv};\addlegendentry{constant sum}
\addplot[omDef,very thick,dotted] table[x=x,y=ss]{figdata/markets-3/20-automated-market-makers/invariants.csv};\addlegendentry{stableswap, $A = 85$}
\end{axis}
\end{tikzpicture}

\omcaption{Three invariants through the balanced point (5, 5), after the \omterm{def:m3:automated-market-makers:stableswap}{stableswap} paper's own comparison: the price is the slope. The
\omterm{def:m3:automated-market-makers:stableswap}{stableswap} curve follows the constant sum over a wide range of balances and bends into a constant-product shape only when one coin
runs short. Curve's integer solver. Data: the chapter's tutorial.}\label{fig:m3:automated-market-makers:invariants}
\end{omfigure}

\section{Impermanent loss and loss-versus-rebalancing}

\begin{definition}[Impermanent loss]\label{def:m3:automated-market-makers:il}
The \emph{impermanent loss}\index{impermanent loss} of a liquidity position is the difference between its value and the value of the
\omterm{def:m3:blockchains-for-traders:key}{tokens} originally deposited, had they simply been held, after a price change; it is ``impermanent'' only in that it vanishes if the
price returns to its starting level.
\end{definition}

\begin{proposition}[Impermanent loss of a constant-product position]\label{prop:m3:automated-market-makers:il}
If the price of the first \omterm{def:m3:blockchains-for-traders:key}{token} moves by a factor $r$, a full-range constant-product position is worth $2\sqrt{r}/(1 + r)$ times what
holding its initial \omterm{def:m3:blockchains-for-traders:key}{tokens} would be worth. A halving or a doubling of the price costs 5.7\%.
\end{proposition}

\begin{proof}
With liquidity $L = \sqrt{xy}$, the pool holds $L/\sqrt{P}$ and $L\sqrt{P}$ and is worth $2L\sqrt{P}$ in the second \omterm{def:m3:blockchains-for-traders:key}{token}. Starting at
$P_0$ and moving to $rP_0$, the pool is worth $2L\sqrt{rP_0}$, and the held \omterm{def:m3:blockchains-for-traders:key}{tokens} $L/\sqrt{P_0} \cdot rP_0 + L\sqrt{P_0} = L\sqrt{P_0}(1 +
r)$. The ratio is $2\sqrt{r}/(1 + r)$.
\end{proof}

\omterm{def:m3:automated-market-makers:il}{Impermanent loss} mixes two things: exposure to the price, which a hedger could remove, and a genuine cost. Milionis, Moallemi,
Roughgarden and Zhang separate them. Compare the pool with a portfolio that holds, at every instant, the same quantity of the risky
\omterm{def:m3:blockchains-for-traders:key}{token} as the pool, bought and sold at market prices as the pool's holding changes: the rebalancing portfolio has the same market risk
as the pool.

\begin{definition}[Loss-versus-rebalancing]\label{def:m3:automated-market-makers:lvr}
\emph{Loss-versus-rebalancing}\index{loss-versus-rebalancing} (LVR) is the difference between the value of a rebalancing portfolio that
tracks an \omterm{def:m3:automated-market-makers:amm}{automated market maker}'s holdings by trading at market prices and the value of the pool itself; it is the amount the pool
loses to arbitrageurs by trading at its own stale prices rather than at the market's.
\end{definition}

\begin{proposition}[LVR of a constant-product pool]\label{prop:m3:automated-market-makers:lvrcp}
If the market price follows a geometric Brownian motion (One Quant Book 4) with volatility $\sigma$ and arbitrageurs keep a fee-free
constant-product pool at the market price, LVR accrues at the instantaneous rate $\sigma^2/8$ per unit of time, as a fraction of the pool's
value. At a daily volatility of 5\%, about 3.1 basis points of pool value a day.
\end{proposition}

\begin{proof}
The paper's general result is an instantaneous LVR of $\tfrac12\sigma^2P^2\,|x^{*\prime}(P)|$, with $x^*(P)$ the pool's holding of the
risky \omterm{def:m3:blockchains-for-traders:key}{token}. For the constant product, $x^* = L/\sqrt{P}$, $|x^{*\prime}| = \tfrac12 L P^{-3/2}$, and the rate is $\tfrac14 L\sigma^2\sqrt{P}$;
divided by the pool value $2L\sqrt{P}$ it is $\sigma^2/8$.
\end{proof}

LVR is the liquidity providers' adverse selection (One Quant Book 1, chapter 19): every arbitrage trade is with someone who knows the price
has moved. It depends on volatility, not on volume. Fees are what the pool earns from everyone else. The tutorial checks the closed form by
simulation: on 200 lognormal paths with 60\% annual volatility, rebalanced every 15 minutes for 30 days, the measured LVR averages 37.1 basis
points of pool value against 37.0 from the formula.

\begin{omfigure}
\begin{tikzpicture}
\begin{axis}[width=11cm,height=5cm,
  xlabel={day},ylabel={cumulative LVR, bp of pool value},
  tick label style={font=\small},label style={font=\small},
  xmin=0,xmax=30,ymin=0,ymax=40,grid=major,grid style={gray!25},
  legend columns=2,legend style={font=\footnotesize,at={(0.5,-0.3)},anchor=north,draw=none,fill=none,/tikz/every even column/.append style={column sep=8pt}},
  table/col sep=comma]
\addplot[omThm,very thick,mark=*,mark size=1pt] table[x=day,y=measured_bp]{figdata/markets-3/20-automated-market-makers/lvr.csv};\addlegendentry{simulated, mean of 200 paths}
\addplot[omDef,thick,dashed] table[x=day,y=theory_bp]{figdata/markets-3/20-automated-market-makers/lvr.csv};\addlegendentry{$\sigma^2/8$ per year}
\end{axis}
\end{tikzpicture}

\omcaption{\omterm{def:m3:automated-market-makers:lvr}{Loss-versus-rebalancing} of a fee-free constant-product pool: the rebalancing portfolio minus the pool, averaged over 200 simulated
paths at 60\% annual volatility with arbitrage every 15 minutes, against the closed form of \cref{prop:m3:automated-market-makers:lvrcp}; the two lines coincide. Data: the
chapter's tutorial.}\label{fig:m3:automated-market-makers:lvr}
\end{omfigure}

\begin{proposition}[When fees pay for LVR]\label{prop:m3:automated-market-makers:breakeven}
If noise traders turn over a volume $V$ a day through a constant-product pool of value $W$ with fee $\gamma_{\mathrm{fee}}$, the pool's fee
income covers its LVR when $V/W \ge \sigma_d^2/(8\gamma_{\mathrm{fee}})$, where $\sigma_d$ is the daily volatility.
\end{proposition}

\begin{proof}
Fee income is $\gamma_{\mathrm{fee}}V$ a day and LVR $\tfrac18\sigma_d^2 W$ a day; compare them.
\end{proof}

At 60\% annual volatility, a 0.30\% pool must turn over 4.1\% of its value every day in non-arbitrage volume just to cover LVR; a 0.05\% pool, 24.7\%.
\Cref{fig:m3:automated-market-makers:turnover} shows the requirement across volatilities. Concentrating liquidity raises fee income and LVR in the same
proportion while the price is in range, so it does not change the ratio; it changes how much capital is at stake and how often the position must be
moved.

\begin{omfigure}
\begin{tikzpicture}
\begin{axis}[width=11cm,height=5cm,ymode=log,log basis y=10,
  xlabel={annual volatility, \%},ylabel={daily volume / pool value, \%},
  tick label style={font=\small},label style={font=\small},
  xmin=20,xmax=120,ymin=0.1,ymax=200,grid=major,grid style={gray!25},
  legend columns=3,legend style={font=\footnotesize,at={(0.5,-0.3)},anchor=north,draw=none,fill=none,/tikz/every even column/.append style={column sep=8pt}},
  table/col sep=comma]
\addplot[omThm,very thick] table[x=vol,y=f5]{figdata/markets-3/20-automated-market-makers/turnover.csv};\addlegendentry{fee 0.05\%}
\addplot[omDef,very thick,dashed] table[x=vol,y=f30]{figdata/markets-3/20-automated-market-makers/turnover.csv};\addlegendentry{fee 0.30\%}
\addplot[omExo,very thick,dotted] table[x=vol,y=f100]{figdata/markets-3/20-automated-market-makers/turnover.csv};\addlegendentry{fee 1\%}
\end{axis}
\end{tikzpicture}

\omcaption{The daily turnover by noise traders at which a constant-product pool's fees equal its \omterm{def:m3:automated-market-makers:lvr}{loss-versus-rebalancing}
(\cref{prop:m3:automated-market-makers:breakeven}), for three fee tiers. Data: the chapter's tutorial.}\label{fig:m3:automated-market-makers:turnover}
\end{omfigure}

\section{Just-in-time liquidity, aggregators and intent-based routing}

\omterm{def:m3:automated-market-makers:concentrated}{Concentrated liquidity} allows a strategy that only a \omterm{def:m3:blockchains-for-traders:chain}{blockchain}'s ordering makes possible.

\begin{definition}[Just-in-time liquidity]\label{def:m3:automated-market-makers:jit}
\emph{Just-in-time liquidity}\index{just-in-time liquidity} is \omterm{def:m3:automated-market-makers:concentrated}{concentrated liquidity} added in a narrow range immediately before a large pending
swap, in the same block, and removed immediately after it, so as to earn most of the swap's fee while bearing almost none of the pool's
arbitrage losses.
\end{definition}

A just-in-time provider sees the swap in the \omterm{def:m3:blockchains-for-traders:mempool}{mempool} (\cref{ch:m3:blockchains-for-traders}), and needs its own transactions ordered before and
after it: it is a form of the ordering games of \cref{ch:m3:maximal-extractable-value}. It takes fees from the pool's passive providers and gives
the trader a better price.

\begin{definition}[DEX aggregator, intent-based routing]\label{def:m3:automated-market-makers:routing}
A \emph{DEX aggregator}\index{DEX aggregator} splits a swap across several pools and venues to minimise the trader's total cost of price impact,
fees and gas. \emph{Intent-based routing}\index{intent-based routing} lets the trader sign a message stating what he will give and the least he will
accept, and lets competing third parties (fillers, solvers) find the execution and submit it on chain.
\end{definition}

\begin{dated}{2026-09}{Two intent-based protocols}\label{dat:m3:automated-market-makers:intents}
UniswapX's documentation describes signed orders that specify a swap's outputs, filled by competing fillers who use their own liquidity or route
through other sources; on some chains the order is a Dutch auction whose price decays from a maximum to a minimum over time, and the swapper pays
no gas and nothing for failed swaps. CoW Protocol's documentation describes signed intents grouped into batches, with third-party solvers competing to
find the best execution for each batch, first by matching opposite intents within it (a ``coincidence of wants'') and otherwise by routing through
AMMs, aggregators and private market makers.
\end{dated}

For a trading firm the consequences are two. As a taker, it can often do better through an aggregator or an intent than against any single pool, and
must compare them net of gas. As a liquidity provider or filler, it competes with specialised firms for the order flow of intents, a business much
closer to the request-for-quote market making of One Quant Book 2 than to passive pool provision.

\section{Tutorial: swaps, impermanent loss and LVR}\label{tut:m3:automated-market-makers:lvr}

\begin{tutorial}
\textbf{Goal.} Implement integer constant-product and \omterm{def:m3:automated-market-makers:stableswap}{stableswap} swaps, compare a liquidity position with holding, and measure \omterm{def:m3:automated-market-makers:lvr}{loss-versus-rebalancing}
by simulation against $\sigma^2/8$. \textbf{End state:} the four figures of this chapter and the numbers of the weekend problem.
\begin{tutsteps}
\item \textbf{The swap.} \Cref{prop:m3:automated-market-makers:swap} in on-chain integer arithmetic.
\omcode{code/firm/amm/firm_amm.py}{14}{22}{The constant-product swap, exact input and exact output.}\label{lst:m3:automated-market-makers:cp}
\item \textbf{LVR.} Hold the pool's risky position in a rebalancing portfolio and compare.
\omcode{code/markets-3/20-automated-market-makers/python/m3_amm.py}{35}{51}{Loss-versus-rebalancing along one simulated path.}\label{lst:m3:automated-market-makers:lvr}
\item \textbf{Run} \texttt{lvr\_mean()}, \texttt{breakeven\_turnover(0.6, 0.003)} and \texttt{fig\_amm.py}.
\end{tutsteps}
\textbf{What to change next.} Give the pool a fee and let arbitrageurs trade only when the mispricing exceeds it; measure how LVR falls with the block time;
simulate a concentrated position that is re-centred when the price leaves its range and count what re-centring costs.
\end{tutorial}

\section{Build: the AMM library}\label{bld:m3:automated-market-makers:amm}

\begin{build}
\bfield{Purpose} The miniature firm must price swaps exactly as the contracts will execute them (to the last unit), value its liquidity positions, and
simulate sandwiches and lending liquidations against pools (\cref{ch:m3:maximal-extractable-value,ch:m3:defi-credit-and-leverage}).
\bfield{Interface} \texttt{cp\_amount\_out}, \texttt{cp\_amount\_in}; \texttt{tick\_at\_price}, \texttt{sqrt\_price\_x96},
\texttt{amounts\_for\_liquidity}, \texttt{next\_sqrt\_price\_from\_amount0}; \texttt{ss\_get\_d}, \texttt{ss\_get\_y}, \texttt{ss\_amount\_out}. The Rust
crate \texttt{firm\_amm} implements the constant-product and \omterm{def:m3:automated-market-makers:stableswap}{stableswap} functions in \texttt{u128}.
\bfield{Rules} Integers only in \omterm{def:m3:blockchains-for-traders:key}{token} base units; floor division where the contracts floor, rounding up where they round up; square-root prices in Q64.96;
Newton iterations capped and raising if they do not converge.
\bfield{Acceptance tests} \texttt{code/firm/amm/tests/} and the Rust crate's tests: the swap formula and the growth of the invariant with fees; ticks and position
amounts at the range edges; the price move after adding a \omterm{def:m3:blockchains-for-traders:key}{token}; \omterm{def:m3:automated-market-makers:stableswap}{stableswap} balance and low slippage against the constant product; one swap with the same
output in Python and Rust.
\bfield{Stretch} Exact tick math in 256-bit integers; multi-range swaps crossing ticks; fees and LVR for a re-centred concentrated position.
\end{build}

\begin{omsources}
\item H.~Adams, N.~Zinsmeister and D.~Robinson, \emph{Uniswap v2 Core}, 2020; H.~Adams et al., \emph{Uniswap v3 Core}, March 2021; Uniswap Labs, ``Uniswap v4 is
here'', 2025.
\item M.~Egorov, \emph{StableSwap --- efficient mechanism for Stablecoin liquidity}, 2019.
\item J.~Milionis, C.~C.~Moallemi, T.~Roughgarden and A.~L.~Zhang, ``Automated Market Making and Loss-Versus-Rebalancing'', arXiv:2208.06046, version of May
2026.
\item UniswapX and CoW Protocol documentation, September 2026; Uniswap Labs, ``Uniswap v3 Mainnet launch'', 5 May 2021; CoinGecko Research, ``DEX to
CEX Volume Ratios'', updated 17 April 2026.
\end{omsources}

\section{Exercises}

\begin{exercise}[$\star$]\label{exo:m3:automated-market-makers:1}
A pool holds 1\,000 ether and 3\,000\,000 dollars with a 0.30\% fee. What does selling it 10 ether return? What is the average price?
\end{exercise}

\begin{exercise}[$\star$]\label{exo:m3:automated-market-makers:2}
What is the \omterm{def:m3:automated-market-makers:il}{impermanent loss} of a full-range position when the price rises by a factor of 4?
\end{exercise}

\begin{exercise}[$\star$]\label{exo:m3:automated-market-makers:3}
Why does a \omterm{def:m3:automated-market-makers:stableswap}{stableswap} pool's liquidity provider effectively sell a put on each coin's peg?
\end{exercise}

\begin{exercise}[$\star\star$]\label{exo:m3:automated-market-makers:4}
Annual volatility is 80\%. What is the LVR of a constant-product pool per day and per year, as a fraction of pool value?
\end{exercise}

\begin{exercise}[$\star\star$]\label{exo:m3:automated-market-makers:5}
A pool of USD~50 million with a 0.05\% fee trades USD~20 million a day of non-arbitrage volume at 50\% annual volatility. Do fees cover LVR?
\end{exercise}

\begin{exercise}[$\star\star$]\label{exo:m3:automated-market-makers:6}
Explain why \omterm{def:m3:automated-market-makers:il}{impermanent loss} and LVR have the same expectation under the risk-neutral measure but are not the same thing.
\end{exercise}

\begin{exercise}[$\star\star\star$]\label{exo:m3:automated-market-makers:7}
\emph{Coding.} With \texttt{lvr\_path}, rebalance every hour instead of every 15 minutes. Does the measured LVR change? What if you rebalance once a day?
\end{exercise}

\begin{exercise}[$\star\star\star$]\label{exo:m3:automated-market-makers:8}
\emph{Find the flaw.} ``Our pool earned 12\% a year in fees, so providing liquidity returned 12\%.''
\end{exercise}

\section{Problem: Should You Provide Liquidity?}

\begin{problem}[{Weekend problem --- fees against loss-versus-rebalancing}]\label{pb:m3:automated-market-makers:1}
A firm considers providing USD~10 million of liquidity to an ether--dollar constant-product pool with a 0.30\% fee. Ether's annual volatility is 60\%.
The pool, USD~200 million in all, trades USD~6 million a day of non-arbitrage volume.

\medskip\noindent\textbf{Part I --- The costs.}
\begin{enumerate}
\item What is the pool's daily LVR as a fraction of its value, and in dollars for the firm's share?
\item What are the firm's daily fees?
\item What daily turnover would make the fees cover LVR?
\item What is the firm's expected annual result, ignoring market risk?
\item What happens to the answer if volatility doubles?
\end{enumerate}

\medskip\noindent\textbf{Part II --- The alternatives.}
\begin{enumerate}[resume]
\item Would the 0.05\% tier with the same volume be better or worse?
\item Would a concentrated $\pm10\%$ range change the ratio of fees to LVR?
\item What does the firm do with its market risk, and what does it cost?
\item What volume would it need at 1\% fee?
\item How does \omterm{def:m3:automated-market-makers:jit}{just-in-time liquidity} change the passive provider's fees?
\end{enumerate}

\medskip\noindent\textbf{Part III --- Measurement.}
\begin{enumerate}[resume]
\item How would you measure the pool's non-arbitrage volume?
\item How would you measure LVR from the pool's history?
\item Why does the block time matter for LVR?
\item What does \omterm{def:m3:automated-market-makers:il}{impermanent loss} over a month tell you about LVR?
\item What would you monitor once the position is live?
\end{enumerate}

\medskip\noindent\textbf{Part IV --- Judgement.}
\begin{enumerate}[resume]
\item Who wins in this pool, and who loses?
\item Is passive liquidity provision a market-making business?
\item Why do pools with high volume and low volatility attract liquidity?
\item State the \emph{named result}: the daily turnover at which fees cover LVR at 60\% volatility and a 0.30\% fee, and whether this pool clears it.
\item In one sentence: what does a liquidity provider sell?
\end{enumerate}
\end{problem}

\section{Interview questions}

\begin{interviewq}[$\star$ \iqroles{trader}]\label{iq:m3:automated-market-makers:1}
Derive the output of a constant-product swap with a fee.
\end{interviewq}

\begin{interviewq}[$\star$ \iqroles{researcher}]\label{iq:m3:automated-market-makers:2}
What is \omterm{def:m3:automated-market-makers:il}{impermanent loss}, and why is it a misleading name?
\end{interviewq}

\begin{interviewq}[$\star\star$ \iqroles{researcher}]\label{iq:m3:automated-market-makers:3}
Derive the LVR of a constant-product pool.
\end{interviewq}

\begin{interviewq}[$\star\star$ \iqroles{developer}]\label{iq:m3:automated-market-makers:4}
Why do AMM contracts use integer arithmetic and square-root prices, and what can go wrong in reimplementing them off chain?
\end{interviewq}

\begin{interviewq}[$\star\star$ \iqroles{trader}]\label{iq:m3:automated-market-makers:5}
How would you decide whether to route a USD~5 million swap through a pool, an aggregator or an intent protocol?
\end{interviewq}

\begin{interviewq}[$\star\star\star$ \iqroles{researcher, developer}]\label{iq:m3:automated-market-makers:6}
Design an AMM that loses less to arbitrageurs. What would you change, and what would it cost?
\end{interviewq}
